\section{Centrality}

\begin{frame}{Key definitions: degree and betweenness centrality}
  \footnotesize
  \begin{columns}[T,onlytextwidth]
    \begin{column}{0.53\textwidth}
      \begin{block}{Degree centrality}
        Local visibility: number of direct neighbours.\par
        \vspace{0.35em}
        Intuition: visible and potentially costly in covert settings.
      \end{block}

      \begin{block}{Betweenness centrality}
        \vspace{-0.45em}
        {\setlength{\abovedisplayskip}{0.25em}%
        \setlength{\belowdisplayskip}{0.3em}%
        \[
        c_B(n)=\sum_{s\neq t\neq n}\frac{\sigma_{st}(n)}{\sigma_{st}}
        \quad
        \hat c_B(n)=\frac{c_B(n)}{(N-1)(N-2)}
        \]
        }

        {\scriptsize $\sigma_{st}$ counts shortest paths from $s$ to $t$;
        $\sigma_{st}(n)$ counts those through $n$;
        $\hat c_B$ normalizes by component size $N$.}
      \end{block}
    \end{column}

    \hfill

    \begin{column}{0.39\textwidth}
      \begin{block}{Visual intuition}
        \centering
        \begin{tikzpicture}[scale=0.62]
          \node[circle,fill=PresTwoPrimary,inner sep=2.5pt,label=above:{\scriptsize hub}] (h) at (0,1.5) {};
          \foreach \a in {20,60,...,340} {
            \node[circle,fill=PresTwoMuted!70,inner sep=1.8pt] (p\a) at ({1.12*cos(\a)},{1.5+1.12*sin(\a)}) {};
            \draw[thin] (h) -- (p\a);
          }
          \node[font=\scriptsize] at (0,0.1) {High degree};

          \node[circle,fill=PresTwoPrimary,inner sep=2.5pt,label=above:{\scriptsize bridge}] (b) at (0,-1.2) {};
          \node[circle,fill=PresTwoMuted!70,inner sep=1.8pt] (l1) at (-1.3,-1.2) {};
          \node[circle,fill=PresTwoMuted!70,inner sep=1.8pt] (l2) at (-2.0,-0.8) {};
          \node[circle,fill=PresTwoMuted!70,inner sep=1.8pt] (l3) at (-2.0,-1.6) {};
          \node[circle,fill=PresTwoMuted!70,inner sep=1.8pt] (r1) at (1.3,-1.2) {};
          \node[circle,fill=PresTwoMuted!70,inner sep=1.8pt] (r2) at (2.0,-0.8) {};
          \node[circle,fill=PresTwoMuted!70,inner sep=1.8pt] (r3) at (2.0,-1.6) {};
          \draw[thin] (b)--(l1) (b)--(r1) (l1)--(l2) (l1)--(l3) (r1)--(r2) (r1)--(r3);
          \node[font=\scriptsize] at (0,-2.25) {Low degree, high betweenness};
        \end{tikzpicture}
      \end{block}
    \end{column}
  \end{columns}

  \note{
    \textbf{Time: 2:30}\par
    Explain degree first, then betweenness exactly as the paper defines it.\par
    State why shortest paths matter in covert networks: extra steps mean extra risk and cost.\par
    Make sure the audience understands why low degree and high betweenness can coexist.
  }
\end{frame}

\begin{frame}{Why betweenness matters in a covert network}
  \small
  \begin{columns}[T,onlytextwidth]
    \begin{column}{0.4\textwidth}
      \begin{block}{What the worked example shows}
        In the supplementary example, the woman-labeled node participates in more shortest paths than any of the four men.
      \end{block}

      \begin{exampleblock}{Interpretation}
        If communication or material transfer follows shortest paths, the high-BC node becomes disproportionately important for brokerage and network stability.
      \end{exampleblock}

      {\footnotesize\color{PresTwoSecondary}
      This is the paper's concrete justification for using betweenness as more than a mathematical abstraction.}
    \end{column}

    \begin{column}{0.6\textwidth}
      \PaperFigureSlot[0.64\textheight]{S1.png}{
        Supplementary figure S1: a worked shortest-path example that makes betweenness concrete.
      }
    \end{column}
  \end{columns}

  \note{
    \textbf{Time: 1:30}\par
    Use this as the paper's own intuition slide for BC.\par
    Emphasize communication, passing items, and brokerage rather than abstract graph theory.\par
    Then move into the data sets and empirical results.
  }
\end{frame}
